%% ma: L11-B16-0093
%% lop: 11
%% chuong: 5
%% bai: 16
%% dang: 11.16.1
%% loai: TL
%% muc: VD
%% dap_an: "-"
%% nguon: "Ảnh giáo viên gửi 25/11/2026: Giới hạn hàm số dạng 0/0 (3 bài, có ĐS), Bài 3 câu 1"
%% kien_thuc: "Lớp 11 Bài 16 (giới hạn của hàm số); hằng đẳng thức x^k-1=(x-1)(x^(k-1)+...+1); tổng 1+2+...+n (Lớp 11 Bài 6 cấp số cộng)"
%% trang_thai: chinh-thuc
%% ghi_chu: "ĐS gốc 8/14 sai, đúng là 18/28 = 9/14."
\begin{ex}
Tính giới hạn $A=\lim\limits_{x\to1}\dfrac{x^{20}-2x+1}{x^{30}-2x+1}$.
\loigiai{$x^{20}-2x+1=(x^{20}-1)-2(x-1)=(x-1)(x^{19}+\dots+x+1-2)$, tương tự mẫu bằng $(x-1)(x^{29}+\dots+x+1-2)$. Vậy $A=\dfrac{20-2}{30-2}=\dfrac9{14}$.}
\end{ex}
